\documentclass[a4paper,10pt]{article}
\usepackage[utf8]{inputenc}
%\usepackage[T2A]{fontenc}
\usepackage[serbian]{babel}
\usepackage{hyperref}

% mathematics
\usepackage{amssymb}
\usepackage{amsthm}

% bibliography
\usepackage[backend=biber]{biblatex}
\addbibresource{refs.bib}

% definitions, theorems, and proofs
\newtheorem{definition}{Definicija}[section]
\newtheorem{theorem}{Teorema}[section]
\newtheorem{corollary}{Posledica}[theorem]
\newtheorem{lemma}[theorem]{Lema}
\theoremstyle{remark}
\newtheorem{example}{Primer}[section]

% misc
\usepackage{enumitem}

%opening
\title{Interaktivni dokazivač u prirodnoj dedukciji za logiku prvog reda}
\author{Stefan Milenković}

\begin{document}

\maketitle

\begin{abstract}
 Ovaj dokument predstavlja uputstvo za izradu seminarskog rada iz predmeta ,,Automatsko rezonovanje'' na master studijama na Matemati\v ckom fakultetu u Beogradu. Ujedno, ovaj dokument predstavlja primer po\v zeljne strukture samog rada. Ovaj paragraf predstavlja sa\v zetak (ili apstrakt) koji treba da postoji u svakom radu.
 Sa\v zetak predstavlja rezime samog rada, tj.
 ~sadr\v zaj rada prepri\v can u nekoliko re\v cenica. Sa\v zetak ne treba da sadr\v zi vi\v se od 100 re\v ci.
\end{abstract}

\section{Uvod}

Moderni automatski dokazivači teorema mogu da dokažu mnoga matematička
tvrđenja postpuno samostalno.
Ipak, ni najuspešniji medju njima još uvek ne mogu da dokažu sve matematičke
teoreme koje su dokazane od strane ljudi.
Pored toga, generisani dokazi nisu uvek lako čitljivi i razumljivi.
Sa druge strane, dokazi koje pišu ljudi su razumljiviji, ali se u njima često
potkradaju greške koje nisu lako uočljive.
Interaktivni dokazivači teorema predstavljaju pristup dokazivanju koji teži da
ima čitkost ljudskih rešenja dok maksimalno koristi računarske resurse.
Naime, interaktivni dokazivači zahtevaju ljudsku interakciju tako što se od
korisnika očekuju da konstantno navodi računar ka ispravnom dokazu.
U medjuvremenu se svaki korak mehanički proverava, a jednostavnija
medjutvrđenja se automatski dokazuju.
Rezultat je potpuno mehanički proveren dokaz koji je često čitljiv ljudima.
Primeri interaktivnih dokazivača koji su upotrebi danas su
Isabelle/HOL \cite{nipkow2002isabelle}, Rocq \cite{rocq2026development} i
Lean \cite{moura2021lean}.

Većina modernih dokazivača rezonuje sa logikama višeg reda, a deduktivni
sistemi na kojima se zasnivaju variraju \cite{maric2015survey}.
U ovom radu se ograničavamo na dokazivače u logici prvog reda u prirodnoj
dedukciji.
Prirodna dedukcija je odabrana kao deduktivni sistem koji teži da ima čitljivost
ljudskih dokaza ali je i pogodan za implementaciju na računaru.
Konkretno, cilj ovog rada je (1) opisivanje interaktivnih dokazivača teorema za
logku prvog reda zasnovanih na prirodnoj dedukciji, i (2) kompletna
implementacija jednog takvog dokazivača u programskom jeziku C++.

Ostatak ovog rada organizovan je na sledeći način.
U poglavlju \ref{sec:osnove} navodimo osnovne definicije, pojmove i notaciju
koja će biti korišćena u ostatku rada.
U poglavlju \ref{sec:opis_metode} dajemo detaljni opis algoritma za
primenjivanje pravila prirodne dedukcije za logiku prvog reda.
Poglavlje \ref{sec:implementacija} sadrži detalje implementacije dokazivača.
Najzad, poglavlje \ref{sec:zakljucak} sadrži zaključna razmatranja i osvrt na
metodu opisanu u radu.

\section{Osnove}
\label{sec:osnove}

U ovom poglavlju dajemo uvodimo i dajemo formalan opis sintakse logike prvog
reda i prirodno dedukcije u njoj.

\subsection{Sintaksa logike prvog reda}

Logika prvog reda (ili predikatska logika) je značajno ekspresivnija od iskazne
logike i proširuje njenu sitaksnu.
Jezik nad kojim se logika prvog reda definiše se sastoji od logičkih dela tj.
logičkih simbola i \textit{signature} \cite{janicic2008mlur}.

\begin{definition}
    Logički deo jezika prvog reda čine skupovi:
    \begin{enumerate}
        \item prebrojiv skup promenljivih $V$
        \item skupa logičkih veznika
        $\{\neg, \land, \lor, \rightarrow, \leftrightarrow\}$
        \item skupa kvantifikatora $\{\forall, \exists\}$
        \item skupa logičkih konstanti $\{\top, \bot\}$
        \item skupa pomoćnih simbola $\{(, ), ,\}$.
    \end{enumerate}
\end{definition}

\begin{definition}
    Signatura (ili jezik) $\mathcal{L} = (\Sigma, \Pi, \mathrm{ar})$ je uređena
    trojka koju određuje skup funkcijskih simbola $\Sigma$, skup relacijskih
    (predikatskih) simbola $\Pi$ i funkcija arnosti
    $\mathrm{ar} : \left(\Sigma \cup \Pi\right) \to \mathbb{N}$.
\end{definition}

Reči nad signaturom $\mathcal{L}$ su bilo koji nizovi karaktera koji se sastoje
iz skupa logičkih simbola i skupova $\Sigma$ i $\Pi$.
Ipak, ograničavaćemo se samo na određeni podskup ovog jezika koji ćemo
koristiti za sintaksu logike prvog reda.

\begin{definition}
    Skup termova nad signaturom $\mathcal{L}$ je najmanji skup za koji važi:
    \begin{enumerate}
        \item sve promenljive iz skupa $V$ su termovi
        \item ako je $f \in \Sigma$ funkcijski simbol arnosti $n$ i
        $t_1, t_2, \ldots, t_n$ termovi, onda je i $f(t_1, t_2, \ldots, t_n)$
        term.
    \end{enumerate}
\end{definition}

\begin{definition}
    Skup atomičkih formula nad signaturom $\mathcal{L}$ je namanji skup za koji
    važi:
    \begin{enumerate}
        \item logičke konstante $\top$ i $\bot$ su atomičke formule
        \item ako je $p \in \Pi$ predikatski simbol arnosti $n$ i
        $t_1, t_2, \ldots, t_n$ termovi, onda je i $p(t_1, t_2, \ldots, t_n)$
        atomička formula.
    \end{enumerate}
\end{definition}

Sada možemo da definišemo skup formula logike prvog reda.

\begin{definition}
    Skup formula logike prvog reda nad signaturom $\mathcal{L}$ je najmanji
    skup za koji važi:
    \begin{enumerate}
        \item svaka atomička formula je formula
        \item ako je $A$ formula, onda je i $\neg A$ formula
        \item ako su $A$ i $B$ formule, onda je su i $A \land B$, $A \lor B$,
        $A \rightarrow B$ i $A \leftrightarrow B$ formule
        \item ako je $A$ formula, a $x$ promenljiva, onda su i $\forall x. A$ i
        $\exists x. A$ formule
    \end{enumerate}
\end{definition}

\begin{example}
    \label{ex:fol-formulae}
    Ako nam je data signatura $\mathcal{L}$ takva da je
    $\Sigma = \{0, \mathrm{S}\}$, $\Pi = \{\mathrm{paran}\}$, a funkcija arnosti
    zadata sa $\mathrm{ar}(\mathrm{S}) = 1$, $\mathrm{ar}(\mathrm{paran}) = 2$ i
    $\mathrm{ar}(0) = 0$, tada su formule
    $$\forall x. \mathrm{paran}(x)
    \rightarrow \mathrm{paran}(\mathrm{S}(\mathrm{S}(x))),$$
    $$\mathrm{paran}(0),$$
    $$\neg \mathrm{paran}(\mathrm{S}(0)),$$
    $$\mathrm{paran}(y) \lor \neg\mathrm{paran}(y)$$
    formule prvog reda nad signaturom $\mathcal{L}$.
\end{example}

Razlikujemo vezana i slobodna pojavljivanja promenljivih u formulama prvog
reda.
Intuitivno, pojavljivanje promenljive je vezano ako je ona pod dejstvom nekog
kvantifikatora, a slobodno inače.

\begin{definition}
    Slobodna pojavljivanje promenljive u formuli prvog reda definišemo na
    sledeći način:
    \begin{enumerate}
        \item sva pojavaljivanja promenljivih u atomičkim formulama su slobodna
        \item sva slobodna pojavljivanja promenljivih u formuli $A$ su slobodna
        i u formuli $\neg A$
        \item sva slobodna pojavljivanja promenljivih u formulama $A$ i $B$ su
        slobodna i u formulama $A \land B$, $A \lor B$, $A \rightarrow B$ i
        $A \leftrightarrow B$.
        \item sva slobodna pojavljivanja promenljivih u $A$ različitih od $x$ su
        slobodna i u $\forall x. A$ i $\exists x. A$.
    \end{enumerate}

    Svako pojavljivanje promenljive koje nije slobodno je vezano.
\end{definition}

\begin{example}
    U formulama iz primera \ref{ex:fol-formulae} promenljiva $x$ se pojavaljuje
    vezano, dok se promenljiva $y$ pojavljuje slobodno.
\end{example}

Još jedan sintaksni pojam koji će nam biti važan jeste zamena (supstitucija)
promenljive termom u formuli prvog reda.

\begin{definition}
    Term $t[x \mapsto t']$ je zamena promenljive $x$ u termu $t$ termom $t'$ ako
    je dobijen tako što se na mesto svakog pojavljivanja promenljive $x$ u termu
    $t$ postavi term $t'$. Preciznije:
    \begin{enumerate}
        \item $x[x \mapsto t'] = t'$
        \item $y[x \mapsto t'] = y$, za $y \neq x$
        \item $f(t_1, t_2, \dots, t_n)[x \mapsto t']
        = f(t_1[x \mapsto t'], t_2[x \mapsto t'], \ldots, t_n[x \mapsto t'])$
    \end{enumerate}
\end{definition}

\begin{definition}
    Formula $F[x \mapsto t]$ je zamena promenljive $x$ u formuli $F$ termom $t$
    i definišemo je na sledeći način:
    \begin{enumerate}
        \item $\top [x \mapsto t] = \top$ i $\bot [x \mapsto t] = \bot$
        \item $p(t_1, t_2, \dots, t_n)[x \mapsto t]
        = p(t_1[x \mapsto t], t_2[x \mapsto t], \ldots, t_n[x \mapsto t])$, gde
        je $p \in \Pi$
        \item $(\neg F)[x \mapsto t] = \neg F[x \mapsto t]$
        \item $(F \otimes G)[x \mapsto t]
        = F[x \mapsto t] \otimes G[x \mapsto t]$, gde je
        $\otimes \in \{\land, \lor, \rightarrow, \leftrightarrow\}$
        \item $(Qx. F)[x \mapsto t] = Qx. F$,
        gde je $Q \in \{\forall, \exists\}$
        \item $(Qy. F)[x \mapsto t] = Qy. F[x \mapsto t]$, gde je
        $Q \in \{\forall, \exists\}$, $y \neq x$ i term $t$ ne sadrži
        pojavljivanja promenljive $y$
        \item $(Qy. F)[x \mapsto t] = Qy'. F[y \mapsto y'][x \mapsto t]$, gde je
        $Q \in \{\forall, \exists\}$, $y \neq x$, term $t$ sadrži pojavljivanja
        promenljive $y$ i promenljiva $y'$ se ne pojavljuje ni u $F$ ni u $t$.
    \end{enumerate}
\end{definition}

\begin{example}
    Primeri zamena nad signaturom $\mathcal{L}$ iz primera \ref{ex:fol-formulae}
    mogu biti:
    \begin{enumerate}[label=(\arabic*)]
        \item $\mathrm{paran}(x)[x \mapsto \mathrm{S}(\mathrm{S}(0))]
        = \mathrm{paran}(\mathrm{S}(\mathrm{S}(0)))$
        \item $(\exists x. \mathrm{paran}(x))[x \mapsto 0]
        = \exists x. \mathrm{paran}(x)$
        \item $(\exists x. \mathrm{paran}(y))[x \mapsto \mathrm{S}(y)]
        = \exists y'. \mathrm{paran}(\mathrm{S}(y))$
    \end{enumerate}
\end{example}

\subsection{Prirodna dedukcija}

Prirodna dedukcija i pravila prirodne dedukcije.

\section{Opis metode}
\label{sec:opis_metode}

Ovo poglavlje treba da sad\v zi detaljniji opis samog algoritma koji je implementiran u okviru seminarskog rada. Algoritam mo\v ze biti opisan re\v cima ili pseudokodom (ne detaljno, ve\' c samo na nivou strukture). Treba poseban akcenat staviti na klju\v cne delove algoritma i njihovu ulogu. Po\v zeljno je navesti i primere koji dodatno razja\v snjavaju rad algoritma.

\section{Implementacija}
\label{sec:implementacija}

U ovom poglavlju treba dati detalje implementacije. Treba navesti u kom programskom jeziku je metoda implementirana, kako je k\^od organizovan, koje su klju\v cne funkcije/klase/metode i koja je njihova uloga.

Takodje je bitno dati upustvo za prevodjenje i pokretanje projekta, kao i navesti softver koji je potreban za prevodjenje i pokretanje programa (prevodilac, alati, dodatne bibiloteke i sl.).

\section{Zaklju\v cak}
\label{sec:zakljucak}



\printbibliography

\end{document}
